\chapter{Weak Maximum Principles for Scalars, Tensors, and Systems}
\label{chow-ii-ch10-weak-maximum-principles}
\dcref{ch10}
\source{chow-et-al-ricci-flow-part-ii:page-0028}

This chapter records weak maximum principles for scalar functions, symmetric
two-tensors, and sections of vector bundles, together with the convex-geometric
tools used in their proofs and several analytic variants.

\section{Scalar and tensor principles}

\begin{definition}[semilinear heat operator]
\label{def:chow-ii-semilinear-heat-operator}
\dcref{ch10:eq-10.1}
\group{Scalar maximum principles}
\source{chow-et-al-ricci-flow-part-ii:page-0030}
Let $M^n$ be a closed manifold equipped with metrics $g(t)$.  Let $X(t)$ be a
time-dependent vector field and let $F:\mathbb R\times[0,T]\to\mathbb R$ be
continuous and locally Lipschitz in its first variable.  Set
\[
 Lu=\partial_tu-\Delta_{g(t)}u-\langle X(t),\nabla u\rangle-F(u,t).
\]
\end{definition}

\begin{theorem}[comparison for semilinear heat equations]
\label{thm:chow-ii-scalar-comparison}
\dcref{ch10:10.1}
\group{Scalar maximum principles}
\source{chow-et-al-ricci-flow-part-ii:page-0030}
\uses{def:chow-ii-semilinear-heat-operator}
If $Lu\le Lv$ on $M\times[0,T]$ and $u(\cdot,0)\le v(\cdot,0)$, then
\(u\le v\) on $M\times[0,T]$.
\end{theorem}
\begin{proof}
\uses{def:chow-ii-semilinear-heat-operator}
Choose $L_0$ larger than the Lipschitz constant of $F$ on the compact range of
$u$ and $v$, and put $w=e^{-L_0t}(u-v)$.  If $w$ had a positive maximum at
$(x_0,t_0)$ with $t_0>0$, then at that point
\(\nabla w=0\), $\Delta w\le0$, and $\partial_tw\ge0$.  The inequality
\(Lu\le Lv\), together with the Lipschitz bound for $F$, gives
\(\partial_tw-\Delta w-\langle X,\nabla w\rangle+(L_0-F_1)w\le0\), where
$F_1\le L_0$; hence $0< (L_0-F_1)w\le0$, a contradiction.  Thus $w\le0$.
\end{proof}

\begin{theorem}[scalar weak maximum principle]
\label{thm:chow-ii-scalar-wmp}
\dcref{ch10:10.2}
\group{Scalar maximum principles}
\source{chow-et-al-ricci-flow-part-ii:page-0030}
\uses{def:chow-ii-semilinear-heat-operator,thm:chow-ii-scalar-comparison}
If $Lu\le0$ and $u(\cdot,0)\le c$, then, while the solution $\phi$ of
\(\phi'=F(\phi,t),\ \phi(0)=c\) exists, one has $u(x,t)\le\phi(t)$.
\end{theorem}
\begin{proof}
\uses{thm:chow-ii-scalar-comparison}
The function $v(x,t)=\phi(t)$ satisfies $Lv=0$ and $u(\cdot,0)\le v(\cdot,0)$.
Apply the comparison theorem.
\end{proof}

\begin{theorem}[scalar upper and lower bounds]
\label{thm:chow-ii-scalar-two-sided}
\dcref{ch10:10.3}
\group{Scalar maximum principles}
\source{chow-et-al-ricci-flow-part-ii:page-0030}
\uses{thm:chow-ii-scalar-wmp}
If $Lu=0$ and $c_1\le u(\cdot,0)\le c_2$, then the ODE solutions
\(\phi_i'=F(\phi_i,t)$, $\phi_i(0)=c_i$, satisfy
\(\phi_1(t)\le u(x,t)\le\phi_2(t)\).
\end{theorem}
\begin{proof}
\uses{thm:chow-ii-scalar-wmp}
Apply the scalar principle to $u$ for the upper bound and to $-u$ with the
reversed comparison for the lower bound.
\end{proof}

\begin{example}[ODE decay estimate]
\label{ex:chow-ii-ode-decay}
\dcref{ch10:10.5}
\group{Scalar maximum principles}
\source{chow-et-al-ricci-flow-part-ii:page-0031}
If $x\ge0$ and $x'\le-C_1x^{1+\alpha}$ whenever $x\ge C_2$, then
\[
x(t)\le\max\{(C_1\alpha t)^{-1/\alpha},C_2\}.
\]
\end{example}
\begin{proof}
On every interval where $x\ge C_2$, differentiation gives
\((x^{-\alpha})'=-\alpha x^{-1-\alpha}x'\ge C_1\alpha\).  Integrating until
the first time $x=C_2$ yields the displayed estimate; after that time the
definition of the first hitting time gives $x\le C_2$.
\end{proof}

\begin{definition}[null-eigenvector condition]
\label{def:chow-ii-null-eigenvector}
\dcref{ch10:10.6}
\group{Tensor maximum principles}
\source{chow-et-al-ricci-flow-part-ii:page-0033}
A map $\beta$ on symmetric two-tensors satisfies the null-eigenvector condition
if $\omega\ge0$ and $\omega(V,\cdot)=0$ imply $\beta(\omega,t)(V,V)\ge0$.
\end{definition}

\begin{theorem}[maximum principle for symmetric two-tensors]
\label{thm:chow-ii-two-tensor-wmp}
\dcref{ch10:10.7}
\group{Tensor maximum principles}
\source{chow-et-al-ricci-flow-part-ii:page-0033}
\uses{def:chow-ii-null-eigenvector}
Let $\alpha(t)$ be a smooth symmetric two-tensor on a closed evolving manifold
with
\[
\partial_t\alpha\ge\Delta\alpha+\langle X,\nabla\alpha\rangle+\beta(\alpha,t),
\]
where $\beta$ is locally Lipschitz and satisfies the null-eigenvector condition.
If $\alpha(0)\ge0$, then $\alpha(t)\ge0$ for all later times.
\end{theorem}
\begin{proof}
\uses{def:chow-ii-null-eigenvector}
Assume positivity fails and let $t_0$ be the first time at which the minimum
eigenvalue reaches zero.  At a minimizing point choose a unit null eigenvector
$V$ and extend it by parallel transport.  The first-time condition gives
$(\partial_t\alpha)(V,V)\le0$, while the spatial second derivative test gives
$(\Delta\alpha)(V,V)\ge0$ and the gradient term vanishes.  The differential
inequality therefore yields $0\ge\beta(\alpha,t_0)(V,V)\ge0$.  Applying the
standard $e^{-\varepsilon t}$ barrier and letting $\varepsilon\downarrow0$
excludes equality at a first crossing, proving preservation.
\end{proof}

\begin{remark}[limitations on nonnegative Ricci curvature]
\label{rem:chow-ii-two-tensor-limitations}
\dcref{ch10:10.8}
\group{Tensor maximum principles}
\source{chow-et-al-ricci-flow-part-ii:page-0034}
Nonnegative Ricci and sectional curvature are not preserved by Ricci flow on
arbitrary closed manifolds in dimensions at least four; the tensor principle
requires its null-eigenvector hypothesis.
\end{remark}

\section{Vector-bundle systems}

\begin{definition}[heat equation for a bundle section]
\label{def:chow-ii-bundle-heat-equation}
\dcref{ch10:10.20--10.25}
\group{Bundle maximum principles}
\source{chow-et-al-ricci-flow-part-ii:page-0036,chow-et-al-ricci-flow-part-ii:page-0038}
Let $\pi:V\to M$ have fixed fiber metric and time-dependent compatible
connections.  The rough Laplacian is $\Delta=\operatorname{tr}_g(\nabla^D\nabla)$,
and a section satisfies
\[
\partial_tu=\Delta u+\nabla_{X(t)}u+F(u,t).
\]
The associated fiberwise ODE is $U'=F_x(U,t)$.
\end{definition}

\begin{definition}[parallel-invariant fiber set]
\label{def:chow-ii-parallel-invariant-set}
\dcref{ch10:10.10}
\group{Bundle maximum principles}
\source{chow-et-al-ricci-flow-part-ii:page-0038}
A subset $\mathcal K\subset V$ is parallel invariant if parallel translation along
every base path maps each fiber section of $\mathcal K$ into $\mathcal K$.
\end{definition}

\begin{lemma}[spectral criterion for parallel invariance]
\label{lem:chow-ii-spectral-invariance}
\dcref{ch10:10.11}
\group{Bundle maximum principles}
\source{chow-et-al-ricci-flow-part-ii:page-0039}
\uses{def:chow-ii-parallel-invariant-set}
For a self-adjoint bundle endomorphism with ordered eigenvalues
$\lambda_1\ge\cdots\ge\lambda_r$, every set of the form
\[
\{u:G(\lambda_1(u),\ldots,\lambda_r(u))\le c\}
\]
is invariant under parallel translation.
\end{lemma}
\begin{proof}
Parallel translation preserves the fiber metric and carries eigensections to
eigensections with the same eigenvalues.  Thus the displayed spectral condition
is unchanged along every parallel lift.
\end{proof}

\begin{corollary}[cone of nonnegative tensors]
\label{cor:chow-ii-nonnegative-tensor-cone}
\dcref{ch10:10.12}
\group{Bundle maximum principles}
\source{chow-et-al-ricci-flow-part-ii:page-0040}
\uses{lem:chow-ii-spectral-invariance}
The cone of self-adjoint tensors with all eigenvalues nonnegative is parallel
invariant.
\end{corollary}
\begin{proof}
\uses{lem:chow-ii-spectral-invariance}
Use $G(\lambda)=-\min_a\lambda_a$ in the spectral criterion.
\end{proof}

\begin{theorem}[Weinberger--Hamilton maximum principle for systems]
\label{thm:chow-ii-system-wmp}
\dcref{ch10:10.13}
\group{Bundle maximum principles}
\source{chow-et-al-ricci-flow-part-ii:page-0040}
\uses{def:chow-ii-bundle-heat-equation,def:chow-ii-parallel-invariant-set}
Let $\mathcal K$ be closed, parallel invariant, and fiberwise convex.  If every
solution of the associated ODE starting in $\mathcal K_x$ remains in
$\mathcal K_x$, then every solution of the bundle heat equation starting in
$\mathcal K$ remains in $\mathcal K$.
\end{theorem}
\begin{proof}
\uses{def:chow-ii-bundle-heat-equation}
Suppose the conclusion fails and let $t_0$ be the first exit time.  In the fiber
at a point of maximal distance from $\mathcal K$, choose the nearest point $v$
and a unit supporting covector $\ell$; parallel extend both radially.  The
spatial maximum property gives $\ell(\Delta u)\le0$ and
$\ell(\nabla_Xu)=0$.  The ODE invariance gives $\ell(F(v,t_0))\le0$; local
Lipschitz continuity gives $\ell(F(u,t_0))\le C\,d(u,\mathcal K)$.  Thus the
maximal distance $s$ satisfies $d^+s/dt\le Cs$ with $s(0)=0$, a contradiction.
\end{proof}

\begin{theorem}[nonnegative curvature operator]
\label{thm:chow-ii-curvature-operator-preservation}
\dcref{ch10:10.14}
\group{Ricci-flow applications}
\source{chow-et-al-ricci-flow-part-ii:page-0043,chow-et-al-ricci-flow-part-ii:page-0044}
\uses{thm:chow-ii-system-wmp,cor:chow-ii-nonnegative-tensor-cone}
For a closed Ricci flow, nonnegative curvature operator is preserved.
\end{theorem}
\begin{proof}
\uses{thm:chow-ii-system-wmp,cor:chow-ii-nonnegative-tensor-cone}
After Uhlenbeck's bundle identification, the curvature operator obeys
\[
\partial_t\Rm=\Delta\Rm+\Rm^2+\Rm^\#.
\]
If $\Rm\ge0$, then both quadratic terms are nonnegative on every null
eigenvector, so the associated ODE preserves the nonnegative cone.  Apply the
system principle.
\end{proof}

\begin{theorem}[time-dependent system maximum principle]
\label{thm:chow-ii-time-dependent-system-wmp}
\dcref{ch10:10.16}
\group{Bundle maximum principles}
\source{chow-et-al-ricci-flow-part-ii:page-0044}
\uses{def:chow-ii-bundle-heat-equation,thm:chow-ii-system-wmp}
Let $\mathcal K(t)$ be closed, fiberwise convex, parallel invariant sets with
closed space-time track.  If each associated ODE carries $\mathcal K_x(t_0)$ into
$\mathcal K_x(t)$ for later times, then the bundle heat equation carries
$\mathcal K(0)$ into $\mathcal K(t)$.
\end{theorem}
\begin{proof}
\uses{thm:chow-ii-system-wmp}
Use the distance $s(t)=\max_xd(u(x,t),\mathcal K_x(t))$.  At a first positive
distance, nearest-point support covectors and radial parallel extension give the
same estimate as in the time-independent case; the ODE hypothesis supplies the
forward tangent direction of the moving set.  Hence $d^+s/dt\le Cs$, while
$s(0)=0$.  The scalar comparison lemma forces $s\equiv0$.
\end{proof}

\begin{theorem}[Hamilton--Ivey estimate]
\label{thm:chow-ii-hamilton-ivey}
\dcref{ch10:10.17}
\group{Ricci-flow applications}
\source{chow-et-al-ricci-flow-part-ii:page-0045,chow-et-al-ricci-flow-part-ii:page-0046}
\uses{thm:chow-ii-time-dependent-system-wmp}
For a normalized Ricci flow on a closed three-manifold with
$\nu(\cdot,0)\ge-1$, every point with $\nu<0$ satisfies
\[
R\ge|\nu|\bigl(\log|\nu|+\log(1+t)-3\bigr).
\]
\end{theorem}
\begin{proof}
\uses{thm:chow-ii-time-dependent-system-wmp}
The set of curvature operators satisfying the scalar lower bound
$R\ge-3/(1+t)$ and, when $\nu\le-1/(1+t)$, the displayed logarithmic bound is
closed, convex, and parallel invariant.  Along the curvature ODE, the first
quantity has derivative at least $\frac23R^2$, while direct differentiation of
\(R/(-\nu)-\log(-\nu)-\log(1+t)\) gives a nonnegative derivative whenever
$\nu\le-1/(1+t)$.  Proposition 10.16 therefore preserves the set, proving the
estimate.
\end{proof}

\begin{corollary}[Hamilton--Ivey estimate without normalization]
\label{cor:chow-ii-hamilton-ivey-scaling}
\dcref{ch10:10.18}
\group{Ricci-flow applications}
\source{chow-et-al-ricci-flow-part-ii:page-0046}
\uses{thm:chow-ii-hamilton-ivey}
If $C=\inf_x\nu(x,0)>0$ in dimension three, then at every point where
$\nu<0$,
\[
R\ge|\nu|\bigl(\log|\nu|+\log(C^{-1}+t)-3\bigr).
\]
\end{corollary}
\begin{proof}
Apply the normalized estimate after parabolic rescaling by $C$ and undo the
scaling of curvature and time.
\end{proof}

\begin{corollary}[almost nonnegative curvature]
\label{cor:chow-ii-almost-nonnegative-curvature}
\dcref{ch10:10.19}
\group{Ricci-flow applications}
\source{chow-et-al-ricci-flow-part-ii:page-0046}
\uses{thm:chow-ii-hamilton-ivey}
For a normalized three-dimensional flow with $\nu(\cdot,0)\ge-1$, there is a
bounded nonincreasing function $\phi$ with $\phi(s)\to0$ as $s\to\infty$ such
that
\[
\Rm\ge-\phi(R)R-3e^6,
\]
and, wherever $R\ge3e^6$, the additive constant may be omitted.
\end{corollary}
\begin{proof}
Set $\psi(s)=s(\log s-3)$ on $(e^3,\infty)$ and let $\psi^{-1}$ be its
inverse.  The Hamilton--Ivey inequality gives the assertion with
$\phi(s)=\psi^{-1}(s)/s$ above $3e^6$ and $\phi=0$ below it.  Monotonicity and
the limit follow by differentiating $\psi$.
\end{proof}

\begin{theorem}[system maximum principle with avoidance]
\label{thm:chow-ii-system-avoidance-wmp}
\dcref{ch10:10.20}
\group{Bundle maximum principles}
\source{chow-et-al-ricci-flow-part-ii:page-0046,chow-et-al-ricci-flow-part-ii:page-0047,chow-et-al-ricci-flow-part-ii:page-0074}
\uses{def:chow-ii-bundle-heat-equation,prop:chow-ii-avoidance-ode}
Let $\mathcal K(t)$ be closed, convex, parallel invariant sets and
$\mathcal A(t)\subset\mathcal K(t)$ closed avoidance sets.  If every associated
ODE starting in $\mathcal K(t_0)$ either stays in $\mathcal K(t)$ or first hits
$\mathcal A(t)$, then every bundle heat solution starting in $\mathcal K(0)$
and avoiding $\mathcal A(t)$ remains in $\mathcal K(t)$.
\end{theorem}
\begin{proof}
\uses{prop:chow-ii-avoidance-ode}
Assume an exit while the PDE solution stays a positive distance from the
avoidance track.  On the compact interval up to the first exit, the nearest-point
support argument for the moving-set principle applies because all supporting
points remain away from $\mathcal A$.  The maximal distance therefore satisfies
$d^+s/dt\le Cs$ and starts at zero, contradicting the Dini barrier.  Thus an
exit can occur only after the avoidance set is reached.
\end{proof}

\begin{theorem}[existence of necklike points]
\label{thm:chow-ii-necklike-points}
\dcref{ch10:10.21}
\group{Ricci-flow applications}
\source{chow-et-al-ricci-flow-part-ii:page-0047,chow-et-al-ricci-flow-part-ii:page-0048,chow-et-al-ricci-flow-part-ii:page-0049,chow-et-al-ricci-flow-part-ii:page-0050}
\uses{thm:chow-ii-system-avoidance-wmp,thm:chow-ii-hamilton-ivey}
Let a finite-time singular Ricci flow on a closed three-manifold have
$\Rm\ge-1$ initially and not converge after rescaling to a positive-curvature
spherical space form.  For every $\tau<T$ and $\delta>0$ there is
$(x,t)\in M\times[\tau,T)$ such that
\[
\mu(\Rm)^2+\nu(\Rm)^2\le\delta\lambda(\Rm)^2,
\qquad
\lambda(\Rm)\ge\frac1{4(T-t)}.
\]
\end{theorem}
\begin{proof}
\uses{thm:chow-ii-hamilton-ivey,thm:chow-ii-system-avoidance-wmp}
Apply the avoidance principle to the Hamilton--Ivey set intersected with
\[
\frac{(T-t)^\varepsilon(\lambda-\nu)}{(R+6)^{1-2\varepsilon}}\le C.
\]
For sufficiently small $\varepsilon=\varepsilon(\delta)$, direct differentiation
along the curvature ODE shows this quantity is nonincreasing whenever the
avoidance inequalities fail.  Hence either the avoidance set is reached, giving
the two displayed inequalities, or the normalized pinching persists.  In the
latter case a blow-up sequence and compactness produce a positive-curvature
spherical limit, contrary to the hypothesis.
\end{proof}

\section{Dini derivatives and spatial maxima}

\begin{definition}[Dini derivatives]
\label{def:chow-ii-dini-derivatives}
\dcref{ch10:10.22}
\group{Dini derivative tools}
\source{chow-et-al-ricci-flow-part-ii:page-0051}
For a function $f$, define $d^+f/dt$, $d^-f/dt$, $d_+f/dt$, and $d_-f/dt$ as
the limsup/liminf of forward/backward difference quotients.
\end{definition}

\begin{lemma}[Dini criterion for an upper bound]
\label{lem:chow-ii-dini-upper-bound}
\dcref{ch10:10.23}
\group{Dini derivative tools}
\source{chow-et-al-ricci-flow-part-ii:page-0051,chow-et-al-ricci-flow-part-ii:page-0052}
\uses{def:chow-ii-dini-derivatives}
If $f$ is right upper and left lower semicontinuous, $f(0)\le c$, and
$d^+f/dt\le0$ whenever $f>c$, then $f\le c$.
\end{lemma}
\begin{proof}
\uses{def:chow-ii-dini-derivatives}
For $\varepsilon>0$ set $f_\varepsilon=f-\varepsilon(1+t)$.  If it first exceeded
$c$, the one-sided difference quotients at the first crossing would give
$d^+f/dt\ge\varepsilon$, contradicting the hypothesis.  Let $\varepsilon\downarrow0$.
\end{proof}

\begin{corollary}[monotonicity and exponential barriers]
\label{cor:chow-ii-dini-monotonicity}
\dcref{ch10:10.25--10.26}
\group{Dini derivative tools}
\source{chow-et-al-ricci-flow-part-ii:page-0052}
\uses{lem:chow-ii-dini-upper-bound}
If $d^+f/dt\le0$ then $f$ is nonincreasing.  If $f(0)\le0$ and
$d^+f/dt\le Cf$ whenever $f>0$, then $f\le0$.
\end{corollary}
\begin{proof}
Apply the lemma to $f-f(t_0)$ for monotonicity, and to $e^{-Ct}f$ for the
exponential assertion.
\end{proof}

\begin{lemma}[Dini derivative of a spatial maximum]
\label{lem:chow-ii-spatial-maximum-dini}
\dcref{ch10:10.29}
\group{Dini derivative tools}
\source{chow-et-al-ricci-flow-part-ii:page-0053,chow-et-al-ricci-flow-part-ii:page-0054}
\uses{def:chow-ii-dini-derivatives}
If $S$ is sequentially compact and $g,\partial_tg$ are continuous, then
$f(t)=\max_Sg(\cdot,t)$ is locally Lipschitz and
\[
\frac{d^+f}{dt}=\max_{g(s,t)=f(t)}\partial_tg(s,t).
\]
\end{lemma}
\begin{proof}
Choose maximizing points at times $t+h_j$ and pass to a convergent subsequence.
The mean-value theorem gives the upper bound by a maximizer at time $t$; fixing
any maximizer at $t$ gives the reverse bound.  Compactness also gives local
uniform Lipschitz bounds.
\end{proof}

\begin{corollary}[weak scalar heat maximum principle]
\label{cor:chow-ii-heat-wmp}
\dcref{ch10:10.31}
\group{Dini derivative tools}
\source{chow-et-al-ricci-flow-part-ii:page-0055}
\uses{lem:chow-ii-spatial-maximum-dini,cor:chow-ii-dini-monotonicity}
If $\partial_tu\le\Delta_{g(t)}u$ on a closed manifold, then
$u(x,t)\le\max_xu(x,0)$.
\end{corollary}
\begin{proof}
At a spatial maximizer, $\Delta u\le0$.  The spatial-maximum lemma gives
$d^+f/dt\le0$ for $f(t)=\max_xu(x,t)$, and the monotonicity corollary applies.
\end{proof}

\begin{theorem}[continuously varying-set principle]
\label{thm:chow-ii-continuous-set-wmp}
\dcref{ch10:10.33}
\group{Bundle maximum principles}
\source{chow-et-al-ricci-flow-part-ii:page-0056,chow-et-al-ricci-flow-part-ii:page-0057,chow-et-al-ricci-flow-part-ii:page-0058,chow-et-al-ricci-flow-part-ii:page-0059}
\uses{lem:chow-ii-spatial-maximum-dini,cor:chow-ii-dini-monotonicity}
If $\mathcal K(t)$ is a continuously varying family of closed, convex,
parallel-invariant fibers and the ODE $U'=F(U)$ preserves it, then
$u_t=\Delta u+F(u)$ preserves $\mathcal K(t)$.
\end{theorem}
\begin{proof}
Let $s(t)=\max_xd(u(x,t),\mathcal K_x(t))$.  The distance is convex and its
nearest-point support covector annihilates the Laplacian contribution by the
second derivative test and parallel invariance.  Comparing the ODE trajectory
from the nearest point with $u$ gives $d^+s/dt\le Cs$.  Since $s(0)=0$, the
exponential barrier yields $s=0$.
\end{proof}

\section{Convex geometry and ODE invariance}

\begin{lemma}[nearest point and separating half-space]
\label{lem:chow-ii-nearest-point-separation}
\dcref{ch10:10.35--10.36}
\group{Convex geometry}
\source{chow-et-al-ricci-flow-part-ii:page-0060,chow-et-al-ricci-flow-part-ii:page-0061}
Every point outside a closed convex set has a unique nearest boundary point,
depending continuously on the point (and on a continuously varying convex set),
and the hyperplane perpendicular to the joining segment separates the point
from the set.
\end{lemma}
\begin{proof}
Existence follows from coercivity of distance on a closed set.  If two nearest
points existed, their midpoint would be strictly closer by the parallelogram
identity.  The same identity plus compactness proves continuity.  For the
nearest point $V$ to $W$, convexity implies $\langle X-V,W-V\rangle\le0$ for
all $X$ in the set, which is the separating half-space.
\end{proof}

\begin{definition}[tangent cone and support function]
\label{def:chow-ii-tangent-support}
\dcref{ch10:10.38,ch10:10.44--10.47}
\group{Convex geometry}
\source{chow-et-al-ricci-flow-part-ii:page-0061,chow-et-al-ricci-flow-part-ii:page-0062}
The tangent cone at a boundary point is the intersection of all containing
closed half-spaces through that point.  A support function at $V$ is a unit
linear functional $\ell$ with $\ell(X)\le\ell(V)$ on the convex set.
\end{definition}

\begin{lemma}[distance as a supremum of supports]
\label{lem:chow-ii-distance-support}
\dcref{ch10:10.50--10.53}
\group{Convex geometry}
\source{chow-et-al-ricci-flow-part-ii:page-0063,chow-et-al-ricci-flow-part-ii:page-0064}
\uses{def:chow-ii-tangent-support,lem:chow-ii-nearest-point-separation}
For a closed convex set $\mathcal J$ and $W\notin\mathcal J$,
\[
d(W,\mathcal J)=\sup_{V,\ell}\ell(W-V),
\]
where $V$ ranges over boundary points and $\ell$ over support functions at $V$.
The supremum is attained; the distance is convex and $C^1$ off $\mathcal J$,
with gradient the unit vector from the nearest point to $W$.
\end{lemma}
\begin{proof}
The separating half-space from the nearest-point lemma gives one support
attaining the distance, while every supporting half-space gives a value no
larger than the distance.  A supremum of affine functions is convex.  Uniqueness
of the nearest point gives the stated gradient, and continuity follows from
continuous dependence of that point.
\end{proof}

\begin{proposition}[ODE criterion for a fixed convex set]
\label{prop:chow-ii-fixed-convex-ode}
\dcref{ch10:10.54}
\group{Convex geometry}
\source{chow-et-al-ricci-flow-part-ii:page-0064,chow-et-al-ricci-flow-part-ii:page-0065,chow-et-al-ricci-flow-part-ii:page-0066}
\uses{lem:chow-ii-distance-support,cor:chow-ii-dini-monotonicity}
The ODE $U'=F(U,t)$ preserves a closed convex set $\mathcal J$ iff
\(\ell(F(V,t))\le0\) for every boundary point $V$ and support function $\ell$
at $V$.
\end{proposition}
\begin{proof}
If the inequality fails, the derivative of $\ell(U(t))$ at the boundary points
outward.  Conversely, for $s(t)=d(U(t),\mathcal J)$, the distance-support formula
and local Lipschitz continuity give $d^+s/dt\le Cs$ whenever $s>0$; since
$s(t_0)=0$, the Dini barrier yields $s=0$.
\end{proof}

\begin{definition}[forward tangent directions]
\label{def:chow-ii-forward-track-direction}
\dcref{ch10:10.58}
\group{Convex geometry}
\source{chow-et-al-ricci-flow-part-ii:page-0066}
For the closed space-time track $\mathcal L=\{(V,t):V\in\mathcal J(t)\}$,
$D_{(V,t)}\mathcal L$ consists of limits of vectors $W_i$ for which
$V+h_iW_i\in\mathcal J(t+h_i)$ for every sequence $h_i\downarrow0$ after passage
to a subsequence.
\end{definition}

\begin{lemma}[convexity of forward directions]
\label{lem:chow-ii-forward-directions-convex}
\dcref{ch10:10.60}
\group{Convex geometry}
\source{chow-et-al-ricci-flow-part-ii:page-0067}
\uses{def:chow-ii-forward-track-direction}
If every $\mathcal J(t)$ is convex, then $D_{(V,t)}\mathcal L$ is convex.
\end{lemma}
\begin{proof}
Take approximating vectors for two directions and form their convex combinations;
convexity of $\mathcal J(t+h_i)$ gives the required approximation.
\end{proof}

\begin{proposition}[ODE criterion for moving convex sets]
\label{prop:chow-ii-moving-convex-ode}
\dcref{ch10:10.61}
\group{Convex geometry}
\source{chow-et-al-ricci-flow-part-ii:page-0067,chow-et-al-ricci-flow-part-ii:page-0068,chow-et-al-ricci-flow-part-ii:page-0069}
\uses{def:chow-ii-forward-track-direction,lem:chow-ii-distance-support,cor:chow-ii-dini-monotonicity}
The ODE preserves every $\mathcal J(t)$ iff
\(F(V,t)\in D_{(V,t)}\mathcal L\) at every boundary point of the closed track.
\end{proposition}
\begin{proof}
Necessity follows by taking difference quotients of an ODE trajectory.  For
sufficiency, if an ODE trajectory exits, use the first exit interval and its
distance to $\mathcal J(t)$.  Closedness gives left lower semicontinuity and right
continuity.  Supporting points and forward directions yield
$d^+s/dt\le Cs$, contradicting $s(t_1)=0$ by the Dini barrier.
\end{proof}

\begin{proposition}[ODE criterion with an avoidance set]
\label{prop:chow-ii-avoidance-ode}
\dcref{ch10:10.64}
\group{Convex geometry}
\source{chow-et-al-ricci-flow-part-ii:page-0069,chow-et-al-ricci-flow-part-ii:page-0070}
\uses{prop:chow-ii-moving-convex-ode}
Let $\mathcal A(t)\subset\mathcal J(t)$ have closed track.  If an ODE starting
in $\mathcal J(t_0)\setminus\mathcal A(t_0)$ either stays in $\mathcal J(t)$ or
hits $\mathcal A(t)$, then it suffices to require the forward-direction condition
only on $\partial\mathcal L\setminus\mathcal A\mathcal L$.
\end{proposition}
\begin{proof}
On an interval where the trajectory avoids $\mathcal A$, closedness gives a
positive distance from its compact image to the avoidance track.  The proof of
the moving-set criterion then applies on a shorter interval, forcing zero
distance unless the trajectory reaches the avoidance set.
\end{proof}

\section{Maximum principles for weak heat solutions}

\begin{proposition}[weighted $L^2$ monotonicity]
\label{prop:chow-ii-weighted-l2}
\dcref{ch10:10.68}
\group{Weak heat solutions}
\source{chow-et-al-ricci-flow-part-ii:page-0075,chow-et-al-ricci-flow-part-ii:page-0076}
If $\partial_tg=-2h$, $\operatorname{tr}_gh\ge-K_1$, and
$u(\partial_tu-\Delta u)\le0$ with $u\partial_\nu u\le0$, then for every
Lipschitz $\xi$ satisfying $\xi_t+\frac12|\nabla\xi|^2\le0$,
\[
e^{-K_1t}\int u^2e^\xi\,d\mu_t
\]
is nonincreasing.
\end{proposition}
\begin{proof}
Differentiate under the integral, use $\partial_td\mu_t=-(\operatorname{tr}_gh)d\mu_t$,
the differential inequality for $\xi$, and integrate by parts.  The result is
\[
-2e^{-K_1t}\int|\nabla(ue^{\xi/2})|^2d\mu_t\le0,
\]
including the nonpositive boundary term.
\end{proof}

\begin{corollary}[$L^2$ uniqueness]
\label{cor:chow-ii-l2-uniqueness}
\dcref{ch10:10.70}
\group{Weak heat solutions}
\source{chow-et-al-ricci-flow-part-ii:page-0076}
\uses{prop:chow-ii-weighted-l2}
For $\xi=0$, the weighted energy is nonincreasing; zero initial $L^2$ norm
forces $u\equiv0$.
\end{corollary}
\begin{proof}
Apply the proposition with $\xi=0$ and use nonnegativity of the integrand.
\end{proof}

\begin{corollary}[exhaustion energy estimate]
\label{cor:chow-ii-exhaustion-energy}
\dcref{ch10:10.72}
\group{Weak heat solutions}
\source{chow-et-al-ricci-flow-part-ii:page-0077}
\uses{prop:chow-ii-weighted-l2}
The weighted energy monotonicity passes to a complete noncompact limit of
compact exhaustions whenever the approximating solutions converge in $C^2$ on
compact subsets.
\end{corollary}
\begin{proof}
Apply the compact estimate on each exhaustion domain, then use $C^2$ convergence
on compact sets and monotone convergence for the nonnegative integrands.
\end{proof}

\begin{proposition}[dual weighted energy]
\label{prop:chow-ii-dual-l2}
\dcref{ch10:10.73}
\group{Weak heat solutions}
\source{chow-et-al-ricci-flow-part-ii:page-0077,chow-et-al-ricci-flow-part-ii:page-0078}
\uses{prop:chow-ii-weighted-l2}
For a subsolution of the adjoint heat equation and a weight satisfying
$-\xi_t+\frac12|\nabla\xi|^2\le0$, the corresponding backward weighted $L^2$ energy
is nondecreasing in forward time.
\end{proposition}
\begin{proof}
Differentiate the backward weight, integrate by parts, and complete the square:
the derivative is $2e^{-K_1(T_0-t)}\int|\nabla(ve^{\xi/2})|^2d\mu_t\ge0$.
\end{proof}

\begin{definition}[weak heat solution]
\label{def:chow-ii-weak-heat-solution}
\dcref{ch10:10.75}
\group{Weak heat solutions}
\source{chow-et-al-ricci-flow-part-ii:page-0081}
A function $u\in V_{2,\mathrm{loc}}$ is a weak subsolution of
$u_t-\Delta_{g(t)}u=f$ if, for every nonnegative compactly supported test
function $\phi$, the integrated-by-parts inequality
\[
\left(\int u\phi\,d\mu\right)(t_1)-\int_0^{t_1}\!\int
(-\mathcal H u\phi+u\phi_t-\nabla u\cdot\nabla\phi)\,d\mu,dt
\le\int_0^{t_1}\!\int f\phi\,d\mu,dt
\]
holds for a.e. $t_1$.  Reverse inequality defines a supersolution; equality a
weak solution.
\end{definition}

\begin{definition}[weak initial-boundary solution]
\label{def:chow-ii-weak-ibv-solution}
\dcref{ch10:10.76}
\group{Weak heat solutions}
\source{chow-et-al-ricci-flow-part-ii:page-0082}
\uses{def:chow-ii-weak-heat-solution}
A function in $V^{1,0}_{2,0}(\Omega_T)$ is a weak solution of the initial-boundary
value problem if the preceding identity holds for all admissible test functions,
with the initial term $\int_\Omega u_0\phi(\cdot,0)\,d\mu$ included.
\end{definition}

\begin{corollary}[energy estimate for weak solutions]
\label{cor:chow-ii-weak-energy-estimate}
\dcref{ch10:10.77}
\group{Weak heat solutions}
\source{chow-et-al-ricci-flow-part-ii:page-0082,chow-et-al-ricci-flow-part-ii:page-0083}
\uses{def:chow-ii-weak-ibv-solution}
For uniformly equivalent metrics, a weak solution satisfies
\[
\|u\|_{V_2}\le C\bigl(\|u_0\|_{L^2}+\|f\|_{L^2}\bigr).
\]
\end{corollary}
\begin{proof}
Test the weak equation by $u$, use uniform metric equivalence and
$2ab\le a^2+b^2$, and integrate the resulting differential inequality.
\end{proof}

\begin{theorem}[weak maximum principle]
\label{thm:chow-ii-weak-solution-wmp}
\dcref{ch10:10.78}
\group{Weak heat solutions}
\source{chow-et-al-ricci-flow-part-ii:page-0083}
\uses{def:chow-ii-weak-ibv-solution}
If $u\in V^{1,0}_2$ satisfies $u_t-\Delta_{g(t)}u\le f$ weakly, with
$f\in L^{n+2p}_{n+2+p}$ and $p>n+2$, then
\[
\operatorname*{ess\,sup}u\le \sup_{\partial_{\rm par}\Omega_T}u_+
 +C\,\operatorname{Vol}(\Omega)^{1/(n+2)-1/p}\|f\|_{L^{n+2p}_{n+2+p}}.
\]
\end{theorem}
\begin{proof}
Test the weak inequality with truncated powers $(u-k)_+^q$, apply the parabolic
Sobolev inequality, and iterate the resulting level-set estimate with the
De Giorgi exponents.  The iteration drives the measure of every superlevel
set above the displayed right-hand side to zero, yielding the essential bound.
\end{proof}

\section{ABP and elementary variants}

\begin{theorem}[Aleksandrov--Bakelman--Pucci estimate]
\label{thm:chow-ii-elliptic-abp}
\dcref{ch10:10.79}
\group{ABP maximum principles}
\source{chow-et-al-ricci-flow-part-ii:page-0084,chow-et-al-ricci-flow-part-ii:page-0085,chow-et-al-ricci-flow-part-ii:page-0086}
If $a^{ij}V_iV_j\ge\lambda|V|^2$ and $a^{ij}u_{ij}=f$ on a bounded smooth
domain, then
\[
\|u\|_{C^0(\Omega)}\le\|\varphi\|_{C^0(\partial\Omega)}+
\frac{\operatorname{diam}(\Omega)}{n\omega_n^{1/n}\lambda}\|f\|_{L^n(\Omega)}.
\]
\end{theorem}
\begin{proof}
Subtract the boundary supremum and consider the upper concave envelope.  Its
gradient image contains the gradient ball of radius $\sup u/\operatorname{diam}(\Omega)$.
The area formula on the contact set and arithmetic--geometric mean give
\(n^n\lambda^n|\det D^2u|\le|f|^n\).  Comparing the two gradient volumes gives
the estimate for $\sup u$; apply it to $-u$ for the lower bound and combine.
\end{proof}

\begin{theorem}[parabolic ABP estimate]
\label{thm:chow-ii-parabolic-abp}
\dcref{ch10:10.80}
\group{ABP maximum principles}
\source{chow-et-al-ricci-flow-part-ii:page-0087}
For a weakly parabolic operator with bounded measurable coefficients satisfying
the stated determinant and drift bounds, every continuous
$W^{2,1}_{n+1,\mathrm{loc}}$ subsolution obeys the parabolic ABP bound
\[
\sup u\le\sup_{\partial_{\rm par}\Omega_T}u_+
 +c(n)d^{n/(n+1)}e^{c(n)d^{-1}\|b/\mathcal D^*\|_{L^{n+1}}^{n+1}}
 \|f^+/\mathcal D^*\|_{L^{n+1}}.
\]
\end{theorem}
\begin{proof}
The parabolic contact set maps under the space-time gradient map onto a set
whose Jacobian is controlled by $(f^+/\mathcal D^*)^{n+1}$.  The determinant
bound and Holder's inequality give the stated power of the diameter; integrating
the drift along contact curves yields the exponential factor.  The boundary
maximum is added by applying the argument to $(u-k)_+$.
\end{proof}

\begin{proposition}[elementary parabolic variant]
\label{prop:chow-ii-elementary-parabolic-wmp}
\dcref{ch10:10.81}
\group{ABP maximum principles}
\source{chow-et-al-ricci-flow-part-ii:page-0088}
If $(\partial_t-a^{ij}\nabla_i\nabla_j)u\le F(u,x,t)$, $u\le k$ on the
parabolic boundary, and $F(z,x,t)\le Az$ for $z\le K$, then
\[
u\le K\quad\text{on }\Omega_{T_1},\qquad
T_1=\min\{T,A^{-1}\log(K/k)\}.
\]
\end{proposition}
\begin{proof}
Let $m(t)=\max_{\Omega_t}u$.  Before the first time $m$ reaches $K$, the scalar
maximum principle applied to $u_t-a^{ij}\nabla_i\nabla_ju\le Au$ gives
$m(t)\le ke^{At}$.  This is $<K$ for $t<T_1$, contradiction.
\end{proof}

\begin{proposition}[$C^0$ estimate for Ricci--DeTurck flow]
\label{prop:chow-ii-detruck-c0-estimate}
\dcref{ch10:10.82}
\group{Ricci-flow applications}
\source{chow-et-al-ricci-flow-part-ii:page-0088,chow-et-al-ricci-flow-part-ii:page-0089,chow-et-al-ricci-flow-part-ii:page-0090,chow-et-al-ricci-flow-part-ii:page-0091,chow-et-al-ricci-flow-part-ii:page-0092}
\uses{prop:chow-ii-elementary-parabolic-wmp}
For a bounded domain in a complete background with uniformly bounded curvature,
the Ricci--DeTurck Dirichlet solution remains uniformly $C^0$-close to the
background metric for a time depending only on the desired error, dimension,
and curvature bound, independently of the domain.
\end{proposition}
\begin{proof}
The DeTurck equation has principal part $g^{pq}\bar\nabla_p\bar\nabla_qg_{ij}$.
Apply the elementary variant to $\psi=\sum_i\lambda_i^m$, where $\lambda_i$ are
the background-relative eigenvalues.  The differentiated equation and
$g\ge(1-\delta_0)\bar g$ make the quadratic gradient terms nonpositive for
suitable $m$ and $\delta_0$, leaving $\psi_t\le g^{pq}\bar\nabla_p\bar\nabla_q\psi+A\psi$.
The boundary value is $\psi=n$; the variant gives $\psi\le2n$ for a uniform
short time, hence $\lambda_i\le(2n)^{1/m}\le1+\delta$.  The lower bound follows
from the preliminary lower estimate, proving the claim.
\end{proof}

\begin{remark}[historical provenance]
\label{rem:chow-ii-ch10-history}
\dcref{ch10:10.83}
\group{ABP maximum principles}
\source{chow-et-al-ricci-flow-part-ii:page-0092}
The system maximum principle originated in work of Stys, Weinberger, and
Hamilton; time-dependent and avoidance-set forms are developed in the cited
references.
\end{remark}
